\section{分布}\label{chap:distributions}\label{chap:fourier-transform}

\subsection{试验函数与分布}\label{sec:test-functions-distributions}

\subsubsection{动机}\label{subsec:distributions-motivation}

微分方程中自然出现的许多问题要求对函数, 导数, 收敛, 积分等作广义定义. 本小节讨论若干这类问题, 下文将对它们给出恰当的回答.

\begin{enumerate}
\item 在 Chapter~\ref{chap:introduction} 中指出, 形如 $u(x,t)=F(x+t)+G(x-t)$ 的任意二次可微函数都是波动方程 $u_{tt}=u_{xx}$ 的解. 显然, 即使 $F$ 和 $G$ 并非二次可微, 似乎也自然应将 $u$ 称为一个"广义"解. 一个自然的问题是, 此时能赋予 $u_{tt}$ 和 $u_{xx}$ 什么意义; 显然, 它们不能是通常意义下的"函数". 对 Chapter~\ref{chap:conservation-laws-shocks} 中讨论过的双曲守恒律的激波解, 也出现同样的问题.
\item 考虑常微分方程初值问题
\begin{equation}\label{eq:impulse-ode-initial-value-problem}
u'(t)=f_\epsilon(t),\ u(0)=0,
\end{equation}
其中
\begin{equation}\label{eq:unit-impulse-approximating-pulse}
f_\epsilon(t)=\begin{cases}\dfrac1\epsilon&1<t<1+\epsilon\\0&\text{otherwise.}\end{cases}
\end{equation}
显然, 解为
\begin{equation}\label{eq:impulse-ode-piecewise-solution}
u(t)=\begin{cases}0&0\leq t\leq1\\\dfrac{t-1}{\epsilon}&1\leq t\leq1+\epsilon\\1&t\geq1+\epsilon.\end{cases}
\end{equation}
注意, 当 $\epsilon\to0$ 时, $u$ 的极限存在; 它是一个阶梯函数. 函数 $f_\epsilon$ 的积分为一; 当 $\epsilon$ 趋于零时, 它的支集位于越来越短的时间区间内. 将 $f_\epsilon$ 的"极限"视为瞬时单位脉冲是自然的. 问题是能赋予这一极限什么意义, 以及微分方程在极限下以何种意义成立. 许多涉及理想化点奇性的物理问题中也有类似的问题: 点电荷的电场, 点光源发出的光, 等等.
\item 在 Chapter~\ref{chap:introduction} 中概述了通过极小化积分 $\int_\Omega|\nabla u|^2\,\mathrm{d}\bm{x}$ 来求解 Dirichlet 问题. 要将这些想法变成严格的理论, 一个基本要素显然是定义一类使该积分有限的函数; 积分的平方根自然地在这一函数空间上定义一个范数. 事实表明, $C^1(\overline\Omega)$ 限制过强; 在该积分所定义的范数下, 它不是完备度量空间. 自然应考虑其完备化; 这将引出一些函数, 它们的 $\nabla u$ 在差商逐点极限这一经典定义的意义下并不存在.
\item Fourier 变换是处理全空间上常系数偏微分方程的自然工具. 然而, 在通常意义下 Fourier 积分存在的函数类限制颇强; 特别地, 这些函数在无穷远处必须可积. 显然, 对不满足这一限制的函数, 例如常函数, 若能有 Fourier 变换的概念, 将很有用.
\end{enumerate}

广义函数背后的想法大致如下: 给定 $\Omega$ 上的连续函数 $f(\bm{x})$, 可以定义线性映射
\begin{equation}\label{eq:function-induced-test-functional}
\phi\mapsto\int_\Omega f(\bm{x})\phi(\bm{x})\,\mathrm{d}\bm{x}
\end{equation}
将某一适当函数类 (称为试验函数) 映入 $\mathbb{R}$. 将看到, 这一映射具有某些连续性性质. 于是将广义函数定义为试验函数上具有相同连续性性质的线性映射.

由于打算利用广义函数研究微分方程, 一个关键问题是: 如何定义这类函数的导数? 答案是: 利用分部积分. 将要求 $\phi$ 在 $\partial\Omega$ 附近为零, 因而导数 $\partial f/\partial x_j$ 可以定义为映射
\begin{equation}\label{eq:distribution-derivative-motivating-pairing}
\phi\mapsto-\int_\Omega f(\bm{x})\dfrac{\partial\phi}{\partial x_j}(\bm{x})\,\mathrm{d}\bm{x}.
\end{equation}
显然, 这一定义不要求 $f$ 在通常意义下可微; 唯一的可微性要求是针对 $\phi$ 的. 因此, 将选取具有很好的光滑性性质的函数作为试验函数.
